\section{From residual structure to a global flux correction}
\label{sec:global-regularization}
Let the corrected field be
\begin{equation}
 \flux_{\mathrm{corr}}=\flux_{\mathrm{raw}}+\Delta\flux,
 \qquad \Delta\flux=\begin{bmatrix}\Delta\psi_d&\Delta\psi_q\end{bmatrix}^{\mathsf T}.
\end{equation}
At a fixed current the torque change is
\begin{equation}
 \Delta M=\kt(\iq\Delta\psi_d-\id\Delta\psi_q),
\end{equation}
and hence
\begin{equation}
 r_{M,\mathrm{corr}}=r_{M,\mathrm{raw}}
 -\kt(\iq\Delta\psi_d-\id\Delta\psi_q).
 \label{eq:corrected-residual}
\end{equation}
Applying this equation independently at every point would reproduce noise and
choose an arbitrary member of a local solution line. The desired object is a
single field over the complete supported current domain. A general discrete
formulation is
\begin{multline}
 \min_{\Delta\flux,\{c_g\}}
 \sum_k w_k\left[r_{M,k}-c_{g(k)}
 -\kt(i_{q,k}\Delta\psi_{d,k}-i_{d,k}\Delta\psi_{q,k})\right]^2\\
 +\lambda_{\mathrm{smooth}}J_{\mathrm{smooth}}(\Delta\flux)
 +\lambda_{\mathrm{change}}J_{\mathrm{change}}(\Delta\flux)
 +J_{\mathrm{prior}}.
 \label{eq:global-objective}
\end{multline}
The data term flattens residual shape rather than forcing zero. Smoothness
couples neighboring operating points; change regularization keeps the raw map
unless torque provides enough evidence to move it. Weights can represent a
declared covariance, but repeat scatter is not automatically independent sensor
uncertainty. Bounds, parity or other physics enter through $J_{\mathrm{prior}}$
and must remain visible assumptions.
This data-fidelity-plus-penalty structure is generalized Tikhonov
regularization; the penalty encodes prior information rather than another
measurement \cite{hohage2020inverse}.

\subsection{Why the numerical experiment uses a co-energy correction}
A particularly coherent variant introduces one scalar surface,
\begin{equation}
 \Delta\coenergy(i_d,i_q),\qquad
 \Delta\psid=\partial_{i_d}\Delta W',\quad
 \Delta\psiq=\partial_{i_q}\Delta W'.
 \label{eq:coenergy-correction}
\end{equation}
Both components then come from one smooth potential and the mixed derivatives
are reciprocal by construction. Curvature regularization acts on the scalar
surface rather than coordinating two unrelated maps. The experiments use the
repository's gauge-fixed degree-four tensor B-spline basis and penalize its
third derivatives plus correction magnitude. No reflection parity is imposed.

This parameterization is used as a physically motivated variant of the general
objective, not as a theorem that the measured terminal-current field is a true
storage gradient. Torque alone does not determine radial co-energy, and a
conservative low-residual fit can still be physically wrong. The companion
integrability paper owns the admissibility tests; the companion reconstruction
paper owns general potential fitting. Here the potential basis supplies a
transparent global regularizer and a reproducible reciprocity prior.

The change and smoothness weights select a representative inside a large
feasible family. Sensitivity across $\lambda$ must therefore accompany any
reported correction. On real data, where true flux is unavailable, $\lambda$
cannot be selected by flux RMSE; the paper reports a fixed conservative example
and its path rather than optimizing against the same torque residual and calling
that validation.
